\thispagestyle{plain}
\begin{center}
{\fontsize{10}{12}\selectfont Manuscrito de recopilación para transferencia académica\par}
\vspace{6mm}
{\fontsize{17}{20.5}\selectfont\bfseries Aritmética genealógica de los números primos y estructura espectral de la función zeta de Riemann\par}
\vspace{7pt}
{\fontsize{12}{14.5}\selectfont Transporte aritmético, momentos espectrales y criterio de positividad de Weil\par}
\vspace{8pt}
{\fontsize{11.5}{14}\selectfont Oumar Haidara Fall\space\textperiodcentered\space Rubén Ramos Balsa\par}
\vspace{4pt}
{\fontsize{8}{10}\selectfont Coautoría sin prelación de contribución\par}
\vspace{4pt}
{\fontsize{8}{10}\selectfont Integración y recopilación documental generadas por ChatGPT - Astra.\par}
\end{center}
\vspace{3pt}
\phantomsection
\addcontentsline{toc}{section}{Resumen}
\begin{center}\bfseries Resumen\end{center}
\begingroup\fontsize{10.5}{12.8}\selectfont
Se demuestra un criterio equivalente de positividad de la forma aritmética de
Weil, en cada ventana fijada, mediante un complemento de Schur de tamaño
finito que conserva el efecto del espacio infinito de detalles.
Su aplicación establece
$\mathscr W(f,f)\geq\frac12\|f\|_2^2$ en el dominio conjunto
de los lectores soportado en $(-1/18,1/18)$, así como en sus
traslaciones y refinamientos interiores. Se construyen además un
lector que recupera la prueba desde la cola gamma y transportes
reversibles que preservan la forma completa, incluidas las
contribuciones primo, gamma y polar. La cota se transmite a colecciones
enteras de pruebas con colas exponenciales.

Estos resultados se desarrollan desde un núcleo formal matemático
denominado Holografía Modular Triádica, en adelante HMT.
\emph{All Possible Paths} (APP) es la estructura aritmética de
caminos sobre el soporte toroidal $(\mathbb Z/9\mathbb Z)^2$,
cuyas traslaciones cardinales forman un grafo de Cayley aditivo;
sus hojas conservan suma y producto con residuo y cociente. TRIT es
la firma orientada $\{-1,0,+1\}$ que distingue los regímenes
locales. El \emph{Topological Phase Kernel} (TPK) compone
selección, transporte, retorno y actualización de memoria sobre el
estado enriquecido. Las realizaciones aritmética y espectral
utilizan la estructura discreta conjunta del continuo producida
por ese desarrollo.

El producto nativo sobre palabras nonádicas tiene forma normal
única; sus aperturas multiplicativas seleccionan los irreducibles
que la evaluación identifica con los números primos. Las
ocupaciones de esos modos generan el producto de Euler y los
momentos primo--potencia. El límite de las trazas de calor cíclicas
y la transformación de Mellin producen la función completada y
su ecuación funcional. La reconstrucción del interior de APP desde
la frontera levantada, las hojas de continuación y las vacancias
precisan qué información conserva el refinamiento.

El doble círculo relaciona la coordenada espectral con la memoria
de retornos. La compresión nonádica determina el balance $8/81$ y
los modos centrales la razón $5/9$, cuyos momentos admiten una
factorización positiva en toda dimensión. Para la forma aritmética,
el control de detalles y términos cruzados proporciona la reducción
de Schur y la coercividad indicada. El criterio global de Weil
queda formulado mediante la identificación conjunta de los momentos
aritméticos con un Gram positivo, incluidos todos los momentos
mixtos y la cobertura del dominio de prueba. Esa identificación
global permanece como hipótesis de los resultados espectrales que
la utilizan: las pruebas de coercividad aquí establecidas no
constituyen una demostración de la hipótesis de Riemann.

\medskip
\noindent\textbf{Palabras clave:} forma de Weil; coercividad;
complemento de Schur; números primos; función zeta de Riemann;
Holografía Modular Triádica; producto nativo; reconstrucción de
frontera; memoria; momentos mixtos.
\par\endgroup
